% ===== 问题二正文 =====
% 问题二正文；三张示意图的 TikZ 宏在 figures/q2_schematics.tex 中定义。
\subsection{问题二模型建立与求解}\label{sec:q2}

首测给出了源的大致方向，距离仍不确定。第二检测点既要保持信号接收，又要使新的测向明显缩小定位区域，因此需要在接收范围内权衡交会几何与移动代价。

本问的基本思路分两步：第一步，利用首测已知的信息，划出一片"无论源在定位区域的哪个位置，到了这里都保证能收到信号"的安全区域；第二步，在安全区域内比较各候选位置的移动距离和定位效果，选出综合代价最小的一个。

\subsubsection{模型准备}\label{sec:q2-1}

\textbf{（1）符号约定。}记频道为 $f$，源的位置为 $p$，首个检测点为 $x_1$，首测示向度为 $\theta_1$，对应的单位方向向量为 $u(\theta_1)=(\cos\theta_1,\sin\theta_1)$。场地为半径 $1800$ m 的圆 $\Omega$，$B(x,r)$ 表示以 $x$ 为圆心、半径 $r$ 的圆盘。接收半径记为 $R_f$，测角误差为 $\pm\varepsilon$，$\varepsilon=1^\circ$。

\textbf{（2）首测后的定位区域。}首测给出了一个示向度 $\theta_1$，告诉我们源大致在哪个方向。但由于存在 $\pm1^\circ$ 的测角误差，源的实际方向只能确定到 $\theta_1-1^\circ$ 和 $\theta_1+1^\circ$ 之间。在平面上，这个方向范围对应一个以 $x_1$ 为顶点、张角 $2^\circ$ 的楔形区域 $W(x_1,\theta_1)$（不等式写法同第一问的式\eqref{eq:q1-1}）。

除了方向信息，首测还隐含了一条距离信息：既然在 $x_1$ 处收到了信号，源一定在接收半径以内。接收半径不超过 $1500$ m，因此 $\|p-x_1\|\le1500$ m。再加上源必须在场地 $\Omega$ 内，得到源 $p$ 必须同时满足三个条件：

\begin{itemize}
\item \textbf{角度条件：}$p$ 落在张角 $2^\circ$ 的楔形 $W(x_1,\theta_1)$ 内；
\item \textbf{距离条件：}$\|p-x_1\|\le1500$ m；
\item \textbf{场地条件：}$p\in\Omega$。
\end{itemize}

三者的交集就是首测后源可能所在的范围，称为定位区域：

\begin{equation}\label{eq:q2-1}
C_f^{(1)}=\Omega\cap B(x_1,1500)\cap W(x_1,\theta_1).
\end{equation}

直观地看，这是一条又窄又长的带状区域：宽度方向只有 $2^\circ$ 的张角，而长度方向可达 $1500$ m。

\textbf{（3）定位区域的多边形近似。}上述定位区域 $C_f^{(1)}$ 的边界含有圆弧（来自 $1500$ m 圆和场地圆），不是纯粹的多边形，后续计算不方便。为此，将两段圆弧分别替换为包在外面的正 $128$ 边形，得到一个纯多边形 $\widehat C_f^{(1)}$。采用"包在外面"而非"嵌在里面"的近似，目的是保证多边形比真实区域略大：真实源一定落在多边形内部，不会因为近似而被漏掉。后文的选点和清除判断均以 $\widehat C_f^{(1)}$ 为准。

另外，为了衡量一个区域的大小，取能把它完全罩住的最小圆，称为包围圆。记其半径为 $\rho(C)$，圆心为 $z(C)$。包围圆半径越小，说明源的位置越确定。

\textbf{（4）第二次测向的几何作用。}首测后的定位区域是一条张角仅 $2^\circ$ 的窄带，方向已经很精确了，不确定的主要是源沿这条窄带有多远。第二次测向的作用，正是从另一个位置观测源，用一条新的视线把这条窄带"截断"，从而把范围从一整条窄带缩小到一小块区域。

截断效果的好坏，关键取决于两条视线之间的夹角——在测量学中称为交会角 $\gamma$。图\ref{fig:q2-cross}直观地对比了两种选择：

\begin{itemize}
\item \textbf{沿示向方向前进再测（图a）：}机器狗沿着首测指出的方向直走，到达新位置后再次测向。由于两个观测点几乎在同一条直线上，两条视线接近平行，第二条视线只是把窄带截短了一点，剩余部分仍然很长；
\item \textbf{横向偏移后再测（图b）：}机器狗在前进的同时向侧面偏移一段距离，造成两条视线之间出现明显的夹角。两条窄带斜着相交，交集只剩下一个紧凑的小四边形，定位区域大幅缩小。
\end{itemize}

\begin{center}
\begin{minipage}{\linewidth}
\centering
\QTwoCrossFigure
\captionof{figure}{两种第二检测点的比较。沿示向方向前进只能截短楔形；横向偏移后两条视线斜交，定位区域才明显缩小。}
\label{fig:q2-cross}
\end{minipage}
\end{center}

这种差别可以用交会角来定量描述。设两次检测点到源的距离分别为 $d_1$、$d_2$，交会角为 $\gamma$。每条楔形在源附近的宽度约为 $2d_i\varepsilon$（$\varepsilon$ 按弧度制），两条窄带斜交后形成的区域，其对角线长度约为

\begin{equation}\label{eq:q2-2}
D\approx\frac{2(d_1+d_2)\,\varepsilon}{\sin\gamma}.
\end{equation}

从这个公式可以读出两条规律：交会角 $\gamma$ 越接近 $90^\circ$，$\sin\gamma$ 越大，$D$ 越小，定位越精确；第二检测点离源越近（$d_2$ 越小），$D$ 也越小。因此，理想的第二检测点应当尽量靠近源，同时与首测视线拉开足够的横向距离，使交会角尽量大。

\textbf{（5）接收范围的限制。}然而横向偏移不能无限增大，因为存在一个硬约束：机器狗只有在源的接收半径以内才能收到信号。如果第二检测点离源太远，到了那里什么也收不到，前面的移动就全白费了。

麻烦的是，源的确切位置和接收半径都是未知的，无法事先算出"到底能偏多远"。已知的只是接收半径不小于 $1000$ m。本文的应对办法是：用一个保守但可靠的策略——要求候选点到定位区域内每一个可能的源位置距离都不超过 $1000$ m。这样，不管源实际在定位区域的哪个位置、接收半径到底有多大，到达候选点后都一定能收到信号。至于哪个候选点最值得去，则通过在若干代表性的假想源位置上模拟测向效果来评估。

\subsubsection{模型的建立}\label{sec:q2-2}

\textbf{（1）决策内容。}需要确定的量只有一个：第二检测点的位置 $x\in\mathbb R^2$。记机器狗当前位置为 $x_0$，当前定位区域为 $\widehat C_f$。在本问中，机器狗刚做完首测还没移动，所以 $x_0=x_1$，$\widehat C_f=\widehat C_f^{(1)}$。选点时能用到的信息只有当前位置和首测确定的定位区域，干扰源的真实位置 $p$ 是未知的，不参与选点计算。

\textbf{（2）约束条件——保证接收。}第二次测向的前提是能收到信号。

源的接收半径虽然未知，但有一个已知的下界 $1000$ m：任何源的接收半径都不小于 $1000$ m。基于此，本文采用一个保守的选点原则：要求候选点 $x$ 到定位区域 $\widehat C_f$ 内每一个可能的源位置的距离都不超过 $1000$ m。这样，不管源的实际位置和实际接收半径如何，到达 $x$ 后都一定在接收范围内。

满足这一条件的所有位置的集合，称为\textbf{可靠接收区}，记作 $Q_f$。

上述条件要求检查到定位区域内"每一个点"的距离，似乎涉及无穷多次计算。但由于定位区域是一个凸多边形，有一条便利的性质可以利用：凸多边形内任何一点，都可以表示为若干顶点的加权平均（凸组合），因此它到候选点的距离，不可能超过所有顶点到候选点的距离中的最大值。反过来说，只要候选点到每个顶点的距离都不超过 $1000$ m，它到区域内任何一点的距离就自动不超过 $1000$ m。所以，检验一个候选点只需要计算它到多边形各顶点的距离：

\begin{equation}\label{eq:q2-3}
Q_f=\bigl\{x:\ \|x-a\|\le1000\ \text{对}\ \widehat C_f\ \text{的每个顶点}\ a\ \text{成立}\bigr\}.
\end{equation}

从几何上看，$Q_f$ 就是以各顶点为圆心、半径 $1000$ m 的若干圆盘的公共部分（交集），因此它是一个凸区域。

此外，选点时还剔除该频道已测过的位置，记为 $\mathcal X_f^{\rm obs}$。这是因为在同一地点重复检测不能消除误差，已有的测向结果可以直接复用，重返已测位置没有收益。

\textbf{（3）目标函数。}约束条件保证了"到了就能收到信号"，但可靠接收区内的候选位置仍然很多，需要一个标准来评判哪个最好。选点的目标是缩短整个任务的总耗时，因此将两部分时间相加作为候选点的代价：

\begin{enumerate}
\item 从当前位置走到候选点要花多少时间（走得越远代价越高）；
\item 到达候选点并测向之后，还需要多长时间才能最终定位和清除源（测向效果越好，后续代价越低）。
\end{enumerate}

第一部分可以直接算出。第二部分的困难在于：测向后定位区域能缩小多少，取决于源的真实位置——而这恰恰是未知的。不同的源位置会产生不同的示向度，因而导致不同的测后区域。

为了在源位置未知的情况下评估测向效果，本文采用以下办法：在定位区域内选取若干有代表性的"假想源位置"，对每个假想位置模拟一次测向过程，看测后区域能缩小到多大，然后对这些假想情形取平均。具体地，选取以下代表位置：

\begin{itemize}
\item 定位区域包围圆的圆心 $z=z(\widehat C_f)$，代表"源在区域中间"的典型情形；
\item 从多边形顶点中均匀选取至多 $4$ 个，各自向圆心方向平移四分之一的距离，记为 $p^{(1)},\dots,p^{(h)}$（$h\le5$）。选靠近顶点的位置，是为了考虑源在区域边缘的情形；向内平移，是避免代表点恰好落在边界上而导致退化。
\end{itemize}

对于每个候选点 $x$ 和每个假想源位置 $p^{(k)}$，模拟测向的过程如下：从 $x$ 看向 $p^{(k)}$ 的方向就是一个"理想的示向度"。将这个方向当作新的测向结果，用第一问的方法在当前定位区域上增加一条楔形约束，从而得到一个更小的新区域，其包围圆半径记为 $r_k(x)$。$r_k(x)$ 越小，说明这次测向对定位的帮助越大。特别地，若 $x$ 与 $p^{(k)}$ 的距离不超过 $5$ m，按题设设备会直接给出"源就在附近"的近场反馈，可以当场清除。

取移动速度 $v=5$ m/s，可靠清除的半径门限 $r_c=19.9$ m（略小于光学清除所需的 $20$ m）。候选点 $x$ 的综合代价定义为

\begin{equation}\label{eq:q2-5}
J(x)=\frac{\|x-x_0\|}{v}
+\frac1h\sum_{k=1}^{h}\left[\frac{\|x-p^{(k)}\|}{v}+P_k(x)\right].
\end{equation}

三项的含义分别是：

\begin{itemize}
\item \textbf{本次移动代价：}$\|x-x_0\|/v$——从当前位置走到候选点所需的时间；
\item \textbf{后续接近代价：}$h^{-1}\sum_k\|x-p^{(k)}\|/v$——测向完成后还要走到源附近去清除，这里用到各假想源位置的平均距离来估计；
\item \textbf{剩余定位代价：}$h^{-1}\sum_k P_k(x)$——如果一次测向后区域还很大、不够精确，后续还需要额外的观测，这部分用罚时来体现。
\end{itemize}

罚时 $P_k(x)$ 的规则是：如果已经非常接近源（$\le5$ m，触发近场反馈），不需要进一步定位，罚时为零；如果测后包围圆半径已经小于清除门限 $r_c=19.9$ m，说明定位已经足够精确，同样不罚。否则，以 $6$ s 作为一次额外检测的基础代价，半径每超出门限 $1$ m 再加 $0.1$ s 来反映可能增加的路程：

\begin{equation}\label{eq:q2-penalty}
P_k(x)=
\begin{cases}
0, & \|x-p^{(k)}\|\le5,\\
0, & \|x-p^{(k)}\|>5,\ r_k(x)\le r_c,\\
6+\dfrac{0.5\,[r_k(x)-r_c]}{v},
& \|x-p^{(k)}\|>5,\ r_k(x)>r_c.
\end{cases}
\end{equation}

上述罚时只是用于比较各候选点优劣的估计量，实际后续耗时取决于真实反馈。换频和检测的固定时间对所有候选点相同，不影响排序，因此未计入代价。

\textbf{（4）模型汇总。}综上，第二检测点的选择模型为

\begin{equation}\label{eq:q2-6}
\begin{aligned}
\min_{x}\quad & J(x)\\
\text{s.t.}\quad & \|x-a\|\le1000,\quad a\ \text{为}\ \widehat C_f\ \text{的每个顶点},\\
& x\notin\mathcal X_f^{\rm obs}.
\end{aligned}
\end{equation}

第一条约束保证到达后一定能收到信号，第二条排除已测过的位置。目标函数综合衡量"走过去要花多久"和"到了之后测向能把定位区域缩到多小"。

需要说明的是，上述假想源位置是等权重取平均的，它们并不是源位置的概率分布，只是用来在各候选点之间做横向比较的工具。因此，模型\eqref{eq:q2-6}给出的是一条实用的选点规则，而非理论上的全局最优解。实际到达第二检测点并测向后，只使用真实的示向度来更新定位区域，假想位置不再参与后续计算。

\subsubsection{模型的求解}\label{sec:q2-3}

模型\eqref{eq:q2-6}的可行域 $Q_f$ 是一片连续区域，而目标函数中涉及多边形求交和最小外接圆的计算，没有简洁的解析解。好在实际并不需要理论上的精确最优点，找到一个足够好的检测位置即可。因此本文采用候选点枚举法：在可行域内预先布设若干有代表性的候选点，逐一计算代价，取代价最小的那个。

\textbf{（1）可靠接收区确实存在可行点。}在枚举之前，先回答一个基本问题：满足式\eqref{eq:q2-3}条件的点真的存在吗？换言之，是否存在一个位置，使得它到定位区域内所有可能源位置的距离都不超过 $1000$ m？

答案是肯定的，并且可以给出一个明确的可行区域。论证如下：首测定位区域被一个以 $x_1$ 为顶点、张角 $2^\circ$、长度不超过 $\widetilde R=1500\sec(\pi/128)$ m 的扇形所包含。如图\ref{fig:q23-1}所示，可以找到一个半径为

\begin{equation}\label{eq:q2-4}
r_0=\frac{\widetilde R}{2\cos\varepsilon}\approx750\ \text{m},\qquad z_0=x_1+r_0\,u(\theta_1)
\end{equation}

的圆，圆心 $z_0$ 位于示向线上（即沿首测方向前进约 $750$ m 处），恰好能把整个扇形罩在里面。这是因为扇形内任一点到 $z_0$ 的距离，最大值只出现在扇形的顶点 $x_1$ 或远端弧上，而这两处到 $z_0$ 的距离都不超过 $r_0$。

有了这个覆盖圆，可行性就很清楚了。以 $z_0$ 为圆心、半径 $1000-r_0\approx250$ m 的圆盘内，取任意一点 $x$，则对区域内任意可能源位置 $p$，由三角不等式：
$$\|x-p\|\le\|x-z_0\|+\|z_0-p\|\le250+750=1000.$$
也就是说，\textbf{沿示向方向前进约 $750$ m，以该处为圆心画一个半径约 $250$ m 的圆，圆内的任何位置都保证能再次收到信号}。场地边界只会缩小定位区域（从而扩大可靠接收区），因此这个结论对任意首测位置都成立。

\begin{figure}[!htbp]
\centering
\QTwoRegionFigure
\caption{首测楔形、覆盖圆与可靠接收区（按实际比例绘制）。蓝色为可靠接收区 $Q_f$ 的一部分，由包含楔形的三角形的三个顶点按式\eqref{eq:q2-3}画出；虚线圆为覆盖整个楔形的圆 $B(z_0,r_0)$；绿色为以 $z_0$ 为圆心、半径约 $250$ m 的圆盘，整体落在 $Q_f$ 内。}
\label{fig:q23-1}
\end{figure}

补充说明：式\eqref{eq:q2-3}只是保证接收的充分条件，可靠接收区之外并非一定收不到信号。首测收到信号本身还隐含着 $R_f\ge\|p-x_1\|$，若一并利用可以进一步扩大某些方向上的可行范围。本文统一采用 $1000$ m 下界，一是计算简便（只需检查有限个顶点），二是后续多次测向时可以复用同一规则持续维持接收保证。

\textbf{（2）候选点的生成。}候选点分为两组，布设方式见图\ref{fig:q2-cand}。

\textbf{第一组：沿前进方向布点。}从当前位置 $x_0$ 到定位区域的包围圆圆心 $z$ 画一条连线，在这条线的中点和终点各设一个基准位置，再从每个基准位置向两侧各偏移 $50$ m 和 $150$ m，加上圆心 $z$ 本身，共 $9$ 个候选点。这一组的作用是保证候选集中总有一些路程较短、直奔区域中心的选择。

\textbf{第二组：沿长轴方向布点。}由式\eqref{eq:q2-2}的分析可知，两条视线的交会角越大、定位效果越好，而交会角取决于第二检测点相对于首测视线的横向偏移。区域最长的方向（长轴）正是最需要被”截断”的方向。取多边形中相距最远的一对顶点，其连线方向即为长轴。将全部顶点投影到长轴上，量出区域沿长轴的长度 $w$。然后在长轴的 $1/4$、$3/8$、$5/8$、$3/4$ 四个分位点上，各取轴上一点以及两侧偏移 $b_*$ 的两点；再在圆心 $z$ 的两侧各取一个偏移 $b_*$ 的点，共 $14$ 个候选点。横向偏移量 $b_*$ 随区域大小自适应调整：

\begin{equation}\label{eq:q2-7}
b_*=\min\{200,\ \max\{20,\ w/8\}\}\ \text{m}.
\end{equation}

这意味着：区域越长，横向偏移越大（取区域长度的八分之一），但不超过 $200$ m，也不小于 $20$ m。偏移之后是否仍在可靠接收区内，还需按式\eqref{eq:q2-3}逐一检验。分位点中特意加入了靠内的 $3/8$ 和 $5/8$，因为当首测楔形很长时，$1/4$ 和 $3/4$ 处横向偏移后的点往往已经超出可靠接收区，靠内的分位点可以保留更多横向选择的余地。

两组候选点合并后，去除重复点、该频道已测过的位置以及不满足式\eqref{eq:q2-3}的点，至多保留 $23$ 个可行候选。

\begin{figure}[!htbp]
\centering
\QTwoCandidateFigure
\caption{两组候选检测点的布设示意。第一组沿当前位置到区域中心的连线布点，第二组沿定位区域的长轴布点；超出可靠接收区的点剔除。}
\label{fig:q2-cand}
\end{figure}

\textbf{（3）代价的计算与选点。}对保留下来的每个候选点按式\eqref{eq:q2-5}计算代价，取代价最小者作为第二检测点。需要注意的是，横向偏移并非越大越好：偏移过大虽然交会角更理想，但移动距离也随之增加，甚至可能超出可靠接收区。最终选出的点是移动代价与定位收益之间的折中。第三问沿用本问”筛选可靠位置、模拟测向效果、按真实反馈更新”的基本思路，但针对多源任务采用不同的候选构造方式。本问的参数设定服务于单次选点，不直接替代第三问的完整算法。

\FloatBarrier
\textbf{（4）求解步骤。}完整的求解流程见图\ref{fig:q2-flow}，各步骤归纳为算法\ref{alg:q2-probe}。

\begin{center}
\begin{minipage}{\linewidth}
\centering
\QTwoFlowFigure
\captionof{figure}{第二检测点选择模型的求解流程。筛选并评价候选点，实际测向后输出第二检测点与更新后的定位区域；第三问中的多次测向与清除过程另行建模。}
\label{fig:q2-flow}
\end{minipage}
\end{center}

\par\medskip
\noindent\begin{minipage}{\linewidth}
\small
\refstepcounter{algorithm}\label{alg:q2-probe}
\textbf{算法\thealgorithm：第二检测点选择模型的求解}\par\smallskip
\noindent\textbf{输入：}当前位置 $x_0$、首测示向度 $\theta_1$、定位区域 $\widehat C_f$、该频道已测过的位置
\begin{enumerate}
\setlength{\itemsep}{2pt}
\setlength{\parskip}{0pt}
\item 计算定位区域 $\widehat C_f$ 的包围圆圆心 $z$ 和长轴方向，据此生成两组候选点；
\item 去除重复点和已测过的位置，逐一检验到多边形各顶点的距离，剔除不在可靠接收区内的点；
\item 在定位区域内选取若干假想源位置，对每个候选点模拟测向、计算代价 $J(x)$；
\item 取代价最小的候选点作为第二检测点；
\item 前往该点实际测向，用真实的示向度更新定位区域，输出第二检测点与更新后的定位区域。若获得近场反馈（距离源不超过 $5$ m），则标记可在此处直接清除。
\end{enumerate}
\end{minipage}
\par\medskip

\Needspace{62mm}
\subsubsection{结果分析}\label{sec:q2-4}

\textbf{（1）实例说明}为了直观展示第二检测点的定位效果，以单源模块运行记录中随机场景 800035 的频道 $20$ 为例。该频道的首测在原点进行，第二检测点由算法\ref{alg:q2-probe}选出。两次测向的数据和定位区域的变化列于表\ref{tab:q2-example}。

\Needspace{38mm}
\begin{center}
\begin{minipage}{\linewidth}
\centering\small
\captionof{table}{频道 20 的两次测向与定位区域的变化}
\label{tab:q2-example}
\vspace{4pt}
\begin{tabular}{@{}lcrr@{}}
\toprule
观测 & 检测点坐标 / m & 示向度 & 定位区域包围圆半径 / m \\
\midrule
原点首测 & $(0.000,\ 0.000)$ & $125.00^\circ$ & $750.226$ \\
第二次测向 & $(-170.646,\ 568.027)$ & $148.14^\circ$ & $65.012$ \\
\bottomrule
\end{tabular}
\end{minipage}
\end{center}

算法选出的第二检测点距原点 $593$ m，其中沿首测方向前进了约 $563$ m、横向偏移了约 $186$ m。它属于第二组候选中长轴 $3/8$ 分位处、横向偏移 $b_*$ 的那个点（此时区域长度 $w\approx1500$ m，由式\eqref{eq:q2-7}得 $b_*\approx187.5$ m）。该点到首测多边形各顶点的最大距离为 $960$ m，满足式\eqref{eq:q2-3}的接收约束，实际运行时也确实再次收到了信号。

从定位效果看，两条示向线的交会角约 $23^\circ$，仅凭这一次新增测向，包围圆半径就从 $750$ m 急剧缩小到 $65$ m，缩小幅度超过 $90\%$。

\textbf{（2）与理论公式的对照。}由两条示向线的交点可以估计源到两个检测点的距离分别约为 $1000$ m 和 $470$ m。将这两个距离和 $23^\circ$ 的交会角代入式\eqref{eq:q2-2}，得到定位区域对角线的理论估计 $D\approx130$ m，与实际重建的包围圆直径基本吻合。这说明模型准备中推导的交会误差公式确实能反映实际的定位精度。

该实例也验证了横向偏移不能一味求大的判断。若把检测点进一步外移到 $3/4$ 分位处，虽然交会角更大，但该处的横向偏移点已经超出可靠接收区，无法保证收到信号。此外，$65$ m 的包围圆半径虽然比首测时大幅缩小，但尚未达到 $19.9$ m 的清除门限，仍需后续观测才能完成定位和清除——两次测向不一定能一步到位。第三问将在完整任务中进一步处理这一问题。

\Needspace{50mm}
\subsubsection{问题二的结论}\label{sec:q2-5}

综合以上分析，问题二的回答如下：

\begin{itemize}
\item \textbf{候选区域。}第二检测点应选在可靠接收区 $Q_f$ 内——即到首测定位区域各顶点距离均不超过 $1000$ m 的位置。无论首测位置在哪里，该区域都包含一个可以明确确定的圆盘：沿示向方向前进约 $750$ m，以该处为圆心、半径约 $250$ m，圆内的任何位置都能保证接收。
\item \textbf{选择策略。}在可靠接收区内，兼顾两类候选位置：一类沿前进方向布设、路程较短；另一类沿长轴方向横向偏移、交会角更大。横向偏移并非越大越好，而是需要与区域尺度相适应，并满足接收约束。具体由模型\eqref{eq:q2-6}在两组候选点中比较移动路程与测后定位收益，选出综合代价最小者。
\end{itemize}

在题设情况下，该策略始终保留真实源位置不被排除，并保证第二检测点一定能接收到反馈信号。在第二检测点完成检测并获得新的示向度后，根据实际返回的示向度重新计算定位区域，达到清除条件后实施清除。由此，本问得到的是一套兼顾接收保证、定位收益与移动代价的选点方法。
